% !TeX program = xelatex

\documentclass{csarticle}

\graphicspath{{./images/}}

\lstloadlanguages{TeX,C++}
\lstset{defaultdialect=[LaTeX]Tex}


\begin{document}

% Front matter ===================================================================

\title{Article Example using csarticle.cls}
\author{Ricardo Sanz}
\reference{CS-XXX}
\version{0.1}
\date{\today}
\URL{http://www.myweb.com/Article.pdf}

\maketitle

\tableofcontents

% Main matter ===================================================================

\section{Introduction} %=========================================================

This brief article is an example of use of the \textbf{CORESENSE \LaTeX\ template for articles}. 

\section{The CSarticle \LaTeX\ template} %========================================

The \code|csarticle.cls| class file is based on the basic \code|article.cls| class of \LaTeX. It redefines some of the macros and defines new ones to create a homogeneous style for CORESENSE internal documentation. This is how it looks like:

The template is provided as a \LaTeX\ class file (\code|.cls|), a style file (\code|.sty|) and a PDF/PNG image with the CORESENSE logo:

\begin{itemize}
\item \code|csarticle.cls| --- the article class file. 
\item \code|cscommon.sty| --- common definitions for CORESENSE \LaTeX\ documents.
\item \code|images| --- images like the project logo or the EU emblem.
\end{itemize}

The template uses the following fonts: Ubuntu Light for main text body; Montserrat for titles\footnote{Google Montserrat is similar to Adobe Futura but can be used freely.}; and {Ubuntu Mono} for code. 

\finalcover

\end{document}


